% =====================================================================
% DeepSpine: artículo en formato IEEE de conferencia (versión en español)
% Plantilla del curso: hoja carta, dos columnas, Times, "TABLA I" encima
% y "Fig. 1" debajo, máximo tres niveles de título (III. / A. / 1)).
% Compilar: pdflatex -> bibtex -> pdflatex -> pdflatex
% (en Overleaf basta con subir la carpeta y compilar main.tex).
% =====================================================================
\documentclass[conference,letterpaper]{IEEEtran}
\usepackage[utf8]{inputenc}
\usepackage[T1]{fontenc}
\usepackage[spanish,es-noshorthands,es-tabla]{babel}
\usepackage{amsmath}
% Las tablas y figuras nunca aparecen antes del punto del texto donde se escriben
\usepackage{flafter}
\usepackage{graphicx,booktabs,array,xcolor}
\usepackage{url}
\def\UrlBreaks{\do\_\do\-\do/\do.}
\usepackage{tikz,pgfplots}
\usetikzlibrary{positioning,arrows.meta}
\pgfplotsset{compat=1.16}
\pgfkeys{/pgf/number format/.cd,1000 sep={\,},use comma}
\graphicspath{{figuras/}}
% La plantilla rotula "TABLA I", "Fig. 1" y "Abstract"
\addto\captionsspanish{\renewcommand{\tablename}{TABLA}\renewcommand{\figurename}{Fig.}\renewcommand{\abstractname}{Abstract}}
\renewcommand{\IEEEkeywordsname}{Palabras clave}

\begin{document}

\title{DeepSpine: clasificación de discos lumbares en resonancia magnética sagital T2}

\author{\IEEEauthorblockN{Jhon Jairo Perdomo Perdomo, Julián Andrés Silva Morales, Juan Manuel Tengonó Lasso}
\IEEEauthorblockA{\textit{Programa de Ingeniería de Software, Facultad de Ingeniería, Universidad Surcolombiana}\\
\textit{Neiva, Colombia}\\
{\small\ttfamily u20221203988@usco.edu.co}\\
{\small\ttfamily u20221203367@usco.edu.co}\\
{\small\ttfamily u20202193898@usco.edu.co}}}

\maketitle

\begin{abstract}
El diagnóstico de la patología del disco intervertebral lumbar en resonancia magnética (RM) es manual y varía entre observadores. Este trabajo presenta DeepSpine, un flujo completo para clasificar discos lumbares en cortes sagitales T2 en tres clases: Normal, Degeneración discal y Protrusión. A partir de un conjunto público de 515 pacientes se construyó un dataset de 5\,123 recortes por disco (ROI) de 399 pacientes, mediante depuración de series, detección de discos con YOLOv8n, preetiquetado con revisión humana de todos los ROI y división por paciente 70/15/15; el aumento de datos, aplicado dentro de cada partición, llevó el conjunto a 35\,007 imágenes. Se compararon cuatro modelos con la misma división y las mismas métricas: YOLOv8s-cls y ResNet-50 con ajuste fino, y SVM y Random Forest sobre características de la ResNet-50 ajustada. En los 768 ROI originales de prueba, la ResNet-50 obtuvo la mayor exactitud (88,9\,\%; IC 95\,\%: 86,5--91,0) y el mayor F1 macro (0,769), aunque las diferencias entre los cuatro modelos caen dentro del margen de error. Todos reconocen bien la Protrusión (\emph{recall} $\geq 0{,}94$) y fallan en la Degeneración (\emph{recall} de 0,40 a 0,50). Los clasificadores clásicos igualan a las redes solo cuando trabajan sobre características ajustadas a la tarea. El límite parece estar en los datos (etiquetas sin radiólogo, pocos casos de Degeneración y un fuerte sesgo por nivel vertebral) más que en la arquitectura.
\end{abstract}

\begin{IEEEkeywords}
resonancia magnética, disco intervertebral lumbar, clasificación de imágenes, aprendizaje profundo, división por paciente, ResNet-50, YOLOv8.
\end{IEEEkeywords}

\section{Introducción}
% I. Introducción
El dolor lumbar es la principal causa de años vividos con discapacidad en el mundo: en 2020 afectó a unos 619 millones de personas y se proyecta que alcance 843 millones de casos en 2050 \cite{gbd2023}. La resonancia magnética (RM) es la prueba de imagen de referencia para estudiar el disco intervertebral. En los cortes sagitales ponderados en T2, la pérdida de señal y de altura del disco indica degeneración, según la escala de Pfirrmann \cite{pfirrmann2001}, y la extensión del material discal más allá de su contorno define abombamientos y hernias, según la nomenclatura de consenso \cite{fardon2014}.

La lectura de estos estudios es manual, lleva tiempo y depende de la experiencia de quien la hace. Incluso con una nomenclatura estandarizada, la concordancia entre radiólogos sobre el contorno discal es solo de aceptable a moderada \cite{arana2011}. Un sistema de apoyo que clasifique de forma automática el estado de cada disco podría reducir el tiempo de lectura y hacer más homogéneos los criterios. La visión artificial ya se ha aplicado a esta tarea con buenos resultados \cite{sari2025,tan2024,sayed2025scirep,wang2025}, pero muchas cifras proceden de conjuntos de prueba pequeños, de tareas binarias o de particiones que no se declaran por paciente, y bajan al evaluar con datos externos \cite{yi2023,nigru2024}. Cuando un mismo paciente aporta imágenes a entrenamiento y a prueba, el modelo puede reconocer al paciente en lugar de la patología, y las métricas se inflan \cite{brunekreef2024,rumala2023}.

Este trabajo presenta DeepSpine, un flujo completo para clasificar discos lumbares (L1-L2 a L5-S1) en cortes sagitales T2 del \emph{Lumbar Spine MRI Dataset} \cite{sudirman2019}. El planteamiento inicial del proyecto definió cuatro clases (Normal, Degeneración, Abombamiento y Hernia) y tres preguntas de investigación:
\begin{enumerate}
  \item[P1.] \textbf{Análisis exploratorio:} ¿cómo se distribuyen los casos de cada clase en el dataset y qué diferencias visuales (brillo, forma o tamaño) permiten distinguir un disco sano de uno patológico?
  \item[P2.] \textbf{Aprendizaje clásico frente a profundo:} ¿pueden algoritmos clásicos, más sencillos y rápidos (SVM o Random Forest), alcanzar un diagnóstico confiable a partir de características de la imagen, o es necesaria una red neuronal profunda?
  \item[P3.] \textbf{Aprendizaje profundo:} ¿qué arquitectura de red convolucional logra la mayor exactitud (con una meta superior al 90\,\%) y la menor tasa de falsos negativos?
\end{enumerate}
Durante la construcción del dataset, Hernia y Abombamiento resultaron inseparables en el plano sagital y se fusionaron en una sola clase, Protrusión (Sección~\ref{sec:anotacion}); por eso las tres preguntas se responden sobre tres clases.

Las contribuciones son tres: (i) un dataset de 5\,123 ROI de 399 pacientes, construido con un proceso trazable y reproducible desde los datos públicos, que declara sus errores y limitaciones; (ii) una comparación de cuatro modelos de familias distintas (YOLOv8s-cls, ResNet-50, SVM y Random Forest) con la misma división por paciente, las mismas métricas y evaluación sobre imágenes originales; y (iii) un análisis del desbalance de clases y del sesgo por nivel vertebral que relativiza la exactitud alcanzada. La Sección~\ref{sec:relacionados} revisa los trabajos previos; la Sección~\ref{sec:metodologia} describe la construcción del dataset, la formulación del problema y los modelos; la Sección~\ref{sec:resultados} presenta y discute los resultados, y la Sección~\ref{sec:conclusiones} responde a las preguntas y expone las limitaciones y el trabajo futuro.


\section{Trabajos relacionados}
\label{sec:relacionados}
% II. Trabajos relacionados
\subsection{Conjuntos de datos de RM lumbar}
El \emph{Lumbar Spine MRI Dataset} reúne estudios de 515 pacientes con dolor lumbar sintomático \cite{sudirman2019}. Sobre ese mismo corpus, Natalia \emph{et al.} seleccionaron los 515 cortes T2 medio-sagitales y predijeron el grado de Pfirrmann con una exactitud media del 88,1\,\% \cite{natalia2024}, y Masood \emph{et al.} publicaron 1\,545 grados de Pfirrmann validados por un radiólogo para L3-L4, L4-L5 y L5-S1 \cite{masood2026}. Fuera de este corpus, SPIDER ofrece 447 series sagitales T1 y T2 de 218 pacientes de cuatro hospitales, anotadas de forma iterativa \cite{vandergraaf2024}, y LumbarDISC reúne 2\,697 pacientes de ocho instituciones de seis países, anotados por radiólogos expertos \cite{richards2026}. Los conjuntos multicéntricos y anotados por expertos son la referencia para validar modelos, pero el de Sudirman sigue siendo uno de los pocos que permiten trabajar con discos individuales en un plazo académico.

\subsection{Clasificación de discos con aprendizaje profundo}
La familia YOLO es la más usada en los trabajos revisados. Liawrungrueang \emph{et al.} detectaron y graduaron la degeneración con YOLOv5 sobre 1\,000 imágenes T2 sagitales de 515 pacientes, con exactitud superior al 95\,\% \cite{liawrungrueang2023}; Prisilla \emph{et al.} compararon YOLOv5x, YOLOv6 y YOLOv7 para detectar hernias con etiquetas de radiólogo \cite{prisilla2023}; Wang \emph{et al.} graduaron degeneración y hernia con un YOLO modificado en T2 sagital y axial, con 420 pacientes internos y 52 externos \cite{wang2025}. Sayed \emph{et al.} usaron YOLOv8 para detectar los discos y recortar un ROI por disco en cortes sagitales (mAP50 del 99,4\,\% para el disco) \cite{sayed2025scirep}, y un flujo en cascada de detector y CNN para identificar discos y detectar hernias en cortes axiales y sagitales \cite{sayed2025imu}. Con redes residuales, Tan \emph{et al.} combinaron ResNet-50 y YOLOv2 para clasificar discos sanos y herniados, aunque con solo 32 pacientes (160 discos) \cite{tan2024}, y Al-kubaisi y Khamiss aplicaron transferencia de aprendizaje a las mismas cuatro clases de este proyecto en 181 pacientes, sobre todo en cortes axiales \cite{alkubaisi2022}. En tomografía computarizada, abultamiento y protrusión fueron el par de clases más confundido \cite{gao2025}.

\subsection{Características profundas con clasificadores clásicos}
Sobre el mismo corpus de Sudirman, Sarı y Atasoy extrajeron características con ResNet-50, las seleccionaron con optimización por lobo gris (GWO) y las clasificaron con SVM, con un 99,42\,\% en una tarea binaria con cortes T1 y validación cruzada de cinco particiones, sin declarar si la división fue por paciente \cite{sari2025}. Khoiwal y McMillan mostraron que clasificadores lineales sobre \emph{embeddings} congelados igualan a modelos entrenados desde cero en retina, mamografía, dermatoscopia y tórax \cite{khoiwal2024}. En datos tabulares de la columna (ángulos óseos), el mejor clasificador fue un SMOTE-RFE-XGBoost \cite{zhang2023}.

\subsection{Vacío de investigación}
Los resultados más altos de la literatura (97--99\,\%) corresponden a tareas binarias o a conjuntos de prueba pequeños, y las cifras internas bajan en validación externa: en un estudio, el F1 sagital pasó de 0,971 a 0,768 y 0,787 en dos centros externos \cite{yi2023}. Pocos trabajos sobre el corpus de Sudirman declaran una división por paciente, reportan métricas por clase con el desbalance real o analizan sesgos como el nivel vertebral. Este trabajo aborda ese vacío: compara cuatro familias de modelos sobre la misma división por paciente, evalúa sobre imágenes originales y reporta métricas por clase y por nivel.


\section{Metodología}
\label{sec:metodologia}
% =====================================================================
% III. Metodología: construcción del dataset (subsecciones A a G).
% Fragmento sin \section; lo sigue secciones/metodologia_modelos.tex (H a J).
% Niveles de título: \section -> III.  \subsection -> A.  \subsubsection -> 1)
% =====================================================================
\providecommand{\pendiente}[1]{\textcolor{red}{\textbf{[PENDIENTE: #1]}}}
\providecommand{\arch}[1]{\path{#1}}

Esta sección describe, en orden, cómo se pasó del conjunto público original a las 35\,007 imágenes de recortes (ROI) con las que se entrenaron y compararon los cuatro modelos (Secciones~\ref{sec:modelos} y \ref{sec:protocolo}). La Fig.~\ref{fig:flujo} resume el flujo y el número de elementos que pasan por cada etapa. Todas las cifras proceden de los archivos versionados del repositorio.


% ---------------------------------------------------------------------
\subsection{Origen de los datos}
\label{sec:origen}
Se usó el \emph{Lumbar Spine MRI Dataset} (Mendeley Data) \cite{sudirman2019}, que reúne estudios de 515 pacientes con dolor lumbar sintomático y 48\,345 cortes de resonancia magnética, y que ya se ha empleado en trabajos previos sobre clasificación de hernia discal \cite{sari2025} y sobre gradación de la degeneración con la escala de Pfirrmann \cite{natalia2024}. Los estudios se adquirieron con un único equipo, un Siemens MAGNETOM Essenza de 1,5~T, y se distribuyen como archivos DICOM con extensión \arch{.IMA}, en una carpeta por paciente. Las carpetas están numeradas de 0001 a 0575; los 60 números que faltan ya faltan en la fuente. Los cortes tienen un tamaño de píxel de 0,729~mm, un grosor de 4~mm y una separación entre centros de 4,8~mm. Los archivos \texttt{.IMA} conservan sus metadatos DICOM, pero el campo \texttt{PatientID} está vacío en todo el corpus y la fuente solo conserva identificadores de estudio y de serie. Por eso, en este trabajo cada carpeta se trata como un paciente: las 515 carpetas equivalen a los 515 pacientes, y toda mención a ``paciente'', incluida la división por paciente, se refiere a la carpeta.

De los 515 pacientes, 414 conservaron cortes útiles tras la depuración (Sección~\ref{sec:depuracion}) y 399 tras las exclusiones de calidad de los ROI (Sección~\ref{sec:roi}). Cada una de las 414 carpetas corresponde a un \texttt{StudyInstanceUID} distinto y no se repite ningún \texttt{StudyTime}; no hay evidencia de pacientes duplicados, aunque con el identificador vacío no puede descartarse del todo.

\begin{figure}[!htb]
\centering
\begin{tikzpicture}[
  box/.style={draw, rounded corners=2pt, align=left, text width=0.56\columnwidth, inner sep=2pt, font=\scriptsize},
  num/.style={align=left, text width=0.34\columnwidth, font=\scriptsize\bfseries},
  arrow/.style={-{Stealth[length=1.5mm]}, thick},
  node distance=3mm]
  \node[box] (a) {1. Conjunto público (Mendeley), archivos \texttt{.IMA}};
  \node[num, right=2mm of a] {515 pacientes, 48\,345 cortes};
  \node[box, below=of a] (b) {2. Filtrado por nombre de serie (T2 sagital)};
  \node[num, right=2mm of b] {8\,794 cortes};
  \node[box, below=of b] (c) {3. 40\,\% central, DICOM$\to$PNG y limpieza manual};
  \node[num, right=2mm of c] {1\,086 cortes, 414 pacientes};
  \node[box, below=of c] (d) {4. Recorte L1--S1 y detección de discos (ROI 128$\times$128)};
  \node[num, right=2mm of d] {5\,385 ROI};
  \node[box, below=of d] (e) {5. Preetiquetado, revisión humana y exclusiones};
  \node[num, right=2mm of e] {5\,123 ROI, 399 pacientes};
  \node[box, below=of e] (f) {6. División por paciente 70/15/15};
  \node[num, right=2mm of f] {279 / 60 / 60 pacientes};
  \node[box, below=of f] (g) {7. Aumento dentro de cada partición};
  \node[num, right=2mm of g] {35\,007 imágenes};
  \foreach \i/\j in {a/b,b/c,c/d,d/e,e/f,f/g} \draw[arrow] (\i) -- (\j);
\end{tikzpicture}
\caption{Flujo de preparación del dataset y número de elementos en cada etapa.}
\label{fig:flujo}
\end{figure}

% ---------------------------------------------------------------------
\subsection{Depuración de la fuente}
\label{sec:depuracion}
Cada paciente trae varias series y solo una parte sirve para este problema. La Tabla~\ref{tab:descartes} lista lo que se descartó y por qué.

\begin{table}[!htb]
\centering\footnotesize
\caption{Series descartadas de cada estudio y motivo}
\label{tab:descartes}
\begin{tabular}{>{\raggedright\arraybackslash}p{0.33\columnwidth}>{\raggedright\arraybackslash}p{0.58\columnwidth}}
\toprule
Serie & Motivo del descarte \\
\midrule
\arch{LOCALIZER_0001} & Localizadores de muy baja resolución para posicionar al paciente; no aportan información anatómica ni de textura. \\
\arch{POSDISP_[4]_T2_TSE_TRA_...} & Vistas sagitales con líneas superpuestas (\emph{overlays}) que marcan los planos transversales y distorsionan la textura. \\
\arch{T1_TSE_SAG_...}, \arch{T1_TSE_TRA_...} & En T1 la grasa es brillante y el líquido oscuro: no permite evaluar la hidratación del disco, base de la escala de Pfirrmann \cite{pfirrmann2001}, y mezclar T1 y T2 altera los rangos de brillo. \\
\arch{T2_TSE_TRA_...} (axial) & Cortes axiales en lotes de 3 a 5 por nivel; exigirían emparejarlos con cada disco. Se deja para una fase futura. \\
\arch{T2_TIRM_SAG_...} & Inversión-recuperación con supresión grasa, de contraste distinto al T2 TSE. Se eliminó en la depuración manual. \\
\bottomrule
\end{tabular}
\end{table}

El proceso se hizo en cuatro pasos, cada uno con su script (carpeta \arch{scripts/preprocesamiento/}); la Tabla~\ref{tab:etapas} da el número de cortes y de pacientes tras cada uno.
\begin{enumerate}
  \item \arch{ingestion.py} copia los archivos cuyo nombre contiene ``T2'' y ``SAG'', salvo localizadores y \arch{POSDISP}. Ese criterio dejó pasar, además de la serie T2 TSE sagital, la serie T2 TIRM sagital de 44 pacientes.
  \item \arch{limpiar_slices.py} conserva el 40\,\% central de los cortes de cada paciente (descarta el 30\,\% inicial y el 30\,\% final, ordenados por nombre de archivo), donde el canal medular y la región L1--S1 son visibles.
  \item \arch{convertidor_dcm_a_png.py} convierte a PNG de 8 bits con una sola escala por serie: percentiles 1 y 99,5 calculados sobre los píxeles de tejido. No se aplicó CLAHE ni otro realce, para no alterar la señal relativa del T2 que sustenta el criterio diagnóstico; la forma de normalizar la intensidad influye en el desempeño del aprendizaje profundo en resonancia, al menos en segmentación \cite{jimenezpartinen2026}, por lo que se mantuvo una sola normalización, por serie, en todo el corpus.
  \item Depuración manual de los PNG: se eliminaron los estudios del segmento cervical o de rodillas, los cortes T2 sagitales oscuros o con aberraciones, y todos los cortes TIRM.
\end{enumerate}

\begin{table}[!htb]
\centering\footnotesize
\caption{Cortes (por secuencia) y pacientes tras cada paso de la depuración}
\label{tab:etapas}
\begin{tabular}{lrrrr}
\toprule
Paso & TSE & TIRM & Total & Pacientes \\
\midrule
Corpus original & -- & -- & 48\,345 & 515 \\
Filtrado por nombre & 8\,094 & 700 & 8\,794 & 515 \\
40\,\% central & 3\,235 & 284 & 3\,519 & 515 \\
Depuración manual & 1\,086 & 0 & 1\,086 & 414 \\
\bottomrule
\end{tabular}
\end{table}

Tras esta depuración quedaron \textbf{1\,086 cortes sagitales T2 TSE de 414 pacientes} (2,6 cortes por paciente en promedio; entre 1 y 6).

\subsubsection{Error en la selección de cortes}
Al revisar la trazabilidad, una vez construido el dataset, se encontró un error en el paso~2. En los 44 pacientes que tenían serie TIRM, el 40\,\% central se calculó sobre la lista conjunta de ambas series ordenada por nombre. Como los archivos TIRM se ordenan antes que los TSE, la ventana conservada abarcó los cortes centrales del TIRM y los primeros cortes del TSE (del 1 al 6 de 15), que son laterales. En la depuración manual se eliminaron los TIRM y, en 36 de esos 44 pacientes, no quedó ningún corte TSE útil: son 36 de los 101 pacientes que se pierden entre 515 y 414. Los otros 8 entraron con el corte 006 de 15, parasagital y no medio-sagital. Ningún corte TIRM llegó al dataset final. Filtrar solo la serie \arch{T2_TSE_SAG} y elegir el corte medio-sagital dentro de esa serie, como hacen otros trabajos sobre este mismo corpus \cite{natalia2024}, habría conservado esos pacientes. Por falta de tiempo no se rehizo el proceso y el error se declara como limitación. Además, 17 pacientes tienen 30 cortes TSE (posiblemente dos series) y podrían estar afectados por el mismo cálculo; no se verificó.

% ---------------------------------------------------------------------
\subsection{Recorte de la región lumbar y de los ROI por disco}
\label{sec:roi}
\subsubsection{Recorte L1--S1}
Con \arch{aplicar_recortes.py} se recortó cada corte en altura para abarcar principalmente L1--S1, replicando por \emph{template matching} unos recortes manuales previos (mediana de la puntuación de coincidencia 0,9996; solo el corte 0547\_008 quedó por debajo de 0,99). El recorte se extendió 60~px hacia arriba y hacia abajo hasta el borde de la imagen; con ello el relleno necesario en L1-L2 bajó de 779 de 1\,080 recortes (72\,\%) a cero.

\subsubsection{Detector de discos}
Detectar los discos antes de clasificarlos es un paso habitual en la literatura, incluida la familia YOLO sobre cortes sagitales \cite{tan2024,sayed2025scirep}. Para localizarlos se entrenó un detector con un conjunto propio en Roboflow (proyecto \emph{Intervertebral Disks}, versión~1): 100 cortes (70 de entrenamiento, 20 de validación y 10 de prueba), con orientación automática y reescalado a 640$\times$640 con bordes negros, sin aumento de datos en Roboflow. Las cinco clases del detector son los niveles L1-L2 a L5-S1. Se ajustó un YOLOv8n de detección de cajas. El modelo que generó los ROI se entrenó durante 100 épocas; una versión de 200 épocas se descartó porque, al inspeccionar sus detecciones, marcaba discos donde no los había. El detector solo ubica el nivel del disco; no clasifica patología.

Aunque Roboflow no aplicó aumento, Ultralytics aplica por defecto un aumento en línea durante el entrenamiento (mosaico, escala, traslación, variación de color y reflexión horizontal con probabilidad 0,5). La reflexión no cambia el nivel de un disco, porque el orden de L1-L2 a L5-S1 es vertical; en el aumento del dataset de clasificación, en cambio, se excluyó (Sección~\ref{sec:aumento}).

El cuaderno con el que se entrenó el modelo usado no guardó sus métricas. Para reportar el desempeño se repitió el entrenamiento con los mismos datos y la misma arquitectura (50 épocas, paciencia 30, lote 16, imagen de 640~px, semilla 42 y el mismo aumento por defecto; mejor época en validación: 49). La Tabla~\ref{tab:detector} resume esa repetición con el umbral de uso (confianza 0,35).

\begin{table*}[t]
\centering\footnotesize
\caption{Desempeño de la repetición del entrenamiento del detector (corrida \texttt{detector\_discos\_v1\_s42}), usada como estimación del modelo que generó los ROI}
\label{tab:detector}
\begin{tabular}{lrrrrrrr}
\toprule
Partición & Cortes & Precisión & Recall & mAP50 & mAP50-95 & Cortes con 5 niveles & IoU medio \\
\midrule
Validación & 20 & 0,922 & 0,964 & 0,962 & 0,663 & 95\,\% & 0,819 \\
Prueba     & 10 & 0,983 & 0,994 & 0,995 & 0,691 & 100\,\% & 0,843 \\
\bottomrule
\end{tabular}
\end{table*}

El nivel más difícil es L1-L2 (en validación: precisión 0,839, recall 0,820 e IoU 0,74); los demás tienen recall de 1,0. Con solo 20 y 10 cortes, un corte mueve varios puntos porcentuales; es una sola corrida con una sola semilla, y la época se eligió con la validación, por lo que sus cifras están algo infladas. Además, describen la repetición y no el modelo usado, aunque ambos comparten datos y arquitectura y difieren en el número de épocas.

\subsubsection{Extracción del ROI}
Con \arch{autocrop_rois.py} se detecta sobre el recorte y se extrae el parche sobre el PNG completo, sumando el desplazamiento. La caja se amplía (margen posterior 0,40, anterior 0,10 y vertical 0,50), se hace cuadrada y se guarda a 128$\times$128~px, una por nivel. Cada recorte queda registrado en \arch{roi_log.csv} con su caja, su ordenada central, su confianza y el relleno por lado. El resultado fueron \textbf{5\,385 ROI}. El margen anterior de L5-S1 se probó en 0,20 frente a 0,10 y la ganancia fue marginal, por lo que se mantuvo 0,10. La Fig.~\ref{fig:rois} muestra los cinco ROI de un corte.

\begin{figure*}[t]
\centering
\includegraphics[width=0.85\textwidth]{rois_0188_008.png}
\caption{Los cinco ROI originales (L1-L2 a L5-S1) del corte 008 del paciente 0188, tal como quedan en el dataset (128$\times$128~px).}
\label{fig:rois}
\end{figure*}

\subsubsection{Exclusiones por calidad}
Se aplicaron reglas objetivas sobre \arch{roi_log.csv} (función \arch{elegibles()} de \arch{scripts/build_dataset.py}): se descarta el corte completo si falta algún nivel, si dos niveles adyacentes se solapan (IoU $>0{,}5$) o si el orden vertical de los centros no crece de L1-L2 a L5-S1; además se excluyen el paciente 0456, el corte 0547\_008 y los ROI con confianza del detector menor que 0,5. En total se excluyeron \textbf{262 ROI}, que afectan a 48 pacientes; 15 pacientes quedaron fuera por completo. Así, los 5\,385 ROI etiquetados quedan en \textbf{5\,123 ROI de 399 pacientes}. De los ROI excluidos, 207 eran Protrusión, 38 Normal y 17 Degeneración.

% ---------------------------------------------------------------------
\subsection{Anotación y retiquetado}
\label{sec:anotacion}
\subsubsection{Clases}
La guía de anotación (\arch{docs/guias/guia_anotacion_4_clases.md}) define cuatro clases nominales: Normal, Degeneración, Abombamiento y Hernia. Degeneración se definió como pérdida de señal T2 del disco o de su altura sin protrusión, en la línea de la escala de Pfirrmann \cite{pfirrmann2001}. Cuando coinciden varias clases, se aplica la prioridad Hernia $>$ Abombamiento $>$ Degeneración $>$ Normal, que solo sirve para desempatar el etiquetado y no es un orden de gravedad. La guía recoge criterios de uso común y no está validada por un radiólogo. Las etiquetas las asignaron estudiantes, sin radiólogo.

\subsubsection{Fusión a tres clases}
En la primera corrida del clasificador semilla, con las cuatro clases (60 imágenes de validación), la exactitud fue 0,65 (IC 95\,\%: 0,52--0,76) y 10 de los 21 errores fueron confusiones Hernia--Abombamiento. Sin un corte axial, la extensión del abultamiento sobre la circunferencia del disco, que es lo que separa ambas entidades en la nomenclatura de consenso \cite{fardon2014}, no puede medirse. Incluso con una nomenclatura estandarizada, la concordancia entre radiólogos sobre el contorno discal es solo de aceptable a moderada \cite{arana2011}, y en tomografía computarizada se ha reportado que abultamiento y protrusión son el par de clases más confundido \cite{gao2025}. Por eso Hernia y Abombamiento se fusionaron en \textbf{Protrusión}, quedando tres clases: Normal, Degeneración discal y Protrusión. La fusión es una decisión posterior a ver los datos (\emph{post hoc}) y se declara así.

\subsubsection{Preetiquetado y revisión}
Etiquetar 5\,385 ROI desde cero no era viable en el plazo, así que se siguió un esquema de preetiquetado con revisión humana, emparentado con el aprendizaje activo \cite{wang2024al} y con la anotación iterativa en la que un modelo propone y una persona corrige \cite{vandergraaf2024}. Primero se etiquetó a mano en Roboflow una semilla balanceada de unos 300 ROI con las cuatro clases originales: 240 de entrenamiento (60 por clase) y 60 de validación (15 por clase). Roboflow reescaló las imágenes a 224$\times$224 y generó tres salidas por imagen de entrenamiento (rotación de $\pm 5^{\circ}$ y brillo de $\pm 15\,\%$). Las etiquetas se refinaron en tres versiones, a medida que se afinaban los criterios morfológicos y se pasaba a tres clases; en la última se retiraron algunas imágenes. La versión~3, usada para el preetiquetado, tiene 719 imágenes de entrenamiento (ya aumentadas) y 84 de validación (24 Degeneración, 30 Normal y 30 Protrusión). Con ella se entrenó un clasificador pequeño, YOLOv8s-cls (entrada 224$\times$224, hasta 60 épocas, paciencia 15, lote 32, semilla 42). En validación alcanzó una exactitud de 0,833 (IC 95\,\%: 0,739--0,898) y un F1 de 0,784 en Degeneración, 0,852 en Normal y 0,857 en Protrusión; con solo 84 imágenes de validación estas cifras son una guía y no una garantía. El preetiquetador fue una herramienta interna: no es ninguno de los cuatro modelos comparados y su única función fue acelerar y ordenar la revisión.

El modelo preetiquetó los 5\,385 ROI con su clase y su confianza (media 0,946; mediana 1,0) y marcó 2\,265 (42,1\,\%) para revisión prioritaria. Las predicciones de Protrusión se enviaron siempre a revisión (2\,039, de las cuales 1\,881 solo por ese motivo), porque la muestra inicial de 70 ROI sugería que era minoritaria y se esperaba que el modelo, entrenado con clases equilibradas, la sobrepredijera. Las otras 226 se marcaron por confianza menor que 0,7, solo o junto con un margen menor que 0,2 entre las dos primeras clases. La revisión mostró lo contrario: el modelo subestimó la Protrusión.

Después, todos los ROI, marcados o no, pasaron por revisión humana en dos grupos, A y B, de 2\,693 y 2\,692 ROI respectivamente. La clase final es siempre la decisión humana. De los 5\,385 ROI, 3\,366 (62,5\,\%) conservaron la clase del modelo y \textbf{2\,019 (37,5\,\%) fueron corregidos}. La Tabla~\ref{tab:reetiquetado} muestra cómo se reorganizó el dataset.

\begin{table}[!htb]
\centering\footnotesize
\setlength{\tabcolsep}{3.5pt}
\caption{Clase propuesta por el preetiquetador frente a la clase final tras la revisión humana (5\,385 ROI)}
\label{tab:reetiquetado}
\begin{tabular}{lrrrrr}
\toprule
 & \multicolumn{3}{c}{Clase final} & & \\
\cmidrule(lr){2-4}
Predicha & Normal & Degen. & Protr. & Total & Corregidas \\
\midrule
Normal       & 1\,086 & 32  & 1\,387 & 2\,505 & 56,6\,\% \\
Degeneración & 51     & 256 & 534    & 841    & 69,6\,\% \\
Protrusión   & 6      & 9   & 2\,024 & 2\,039 & 0,7\,\% \\
\midrule
Total final  & 1\,143 & 297 & 3\,945 & 5\,385 & 37,5\,\% \\
\bottomrule
\end{tabular}
\end{table}

De los 3\,945 ROI de Protrusión finales, 1\,921 (48,7\,\%) llegaron como Normal o Degeneración, mientras que casi nada de lo que el modelo predijo como Protrusión cambió (15 de 2\,039). La revisión humana, y no el modelo, fijó la distribución de clases del dataset, y esa distribución está muy desbalanceada (Sección~\ref{sec:balance}).

\subsubsection{Fiabilidad de las etiquetas}
Cada ROI fue revisado por una sola persona (A o B); 408 de los 414 pacientes tienen ROI revisados por ambos, pero no hay ROI doble. En consecuencia \textbf{no puede calcularse el coeficiente kappa entre anotadores}. Sí se observan diferencias de criterio: A etiquetó Degeneración en el 6,9\,\% de sus ROI y B en el 4,1\,\%, y corrigió el 35,3\,\% frente al 39,7\,\% de B. No es posible separar cuánto de esa diferencia es criterio y cuánto es composición de la muestra. Las etiquetas de no expertos introducen ruido que reduce la capacidad de generalizar de las redes \cite{shi2024}. Además, como los revisores veían la sugerencia del modelo, cabe un sesgo de anclaje: en radiología se ha observado que un consejo incorrecto de una IA arrastra el diagnóstico de los médicos \cite{prinster2024}.

% ---------------------------------------------------------------------
\subsection{División por paciente}
\label{sec:split}
Cada paciente aporta unos 13 ROI (5\,123/399) de cortes vecinos casi idénticos. Si se dividiera por imagen, un paciente quedaría repartido entre entrenamiento y prueba y el modelo podría reconocer al paciente en lugar de la patología, inflando las métricas (fuga de datos) \cite{brunekreef2024,rumala2023}. Por eso \textbf{cada paciente va entero a una sola partición}.

El reparto es 70/15/15 por pacientes, con semilla 42. El script prueba 20\,000 repartos aleatorios y conserva el que mejor mantiene, en las tres particiones, la proporción de cada clase (con peso triple) y de cada nivel. Se eligieron tres particiones porque la validación sirve para decidir (época, hiperparámetros, umbrales) y la prueba para medir una sola vez al final, con todo ya decidido. La Tabla~\ref{tab:split} muestra el resultado. La intersección de pacientes entre particiones es vacía y se comprueba con aserciones en el script. Los pacientes con al menos un ROI de Degeneración son 81, 17 y 22 en entrenamiento, validación y prueba.

\begin{table}[!htb]
\centering\footnotesize
\setlength{\tabcolsep}{3.5pt}
\caption{División por paciente y ROI originales por clase}
\label{tab:split}
\begin{tabular}{lrrrrrr}
\toprule
Partición & Pacientes & Normal & Degen. & Protr. & ROI & \% ROI \\
\midrule
Entrenamiento & 279 & 771   & 194 & 2\,612 & 3\,577 & 69,8 \\
Validación    & 60  & 170   & 44  & 564    & 778    & 15,2 \\
Prueba        & 60  & 164   & 42  & 562    & 768    & 15,0 \\
\midrule
Total         & 399 & 1\,105 & 280 & 3\,738 & 5\,123 & 100 \\
\bottomrule
\end{tabular}
\end{table}

% ---------------------------------------------------------------------
\subsection{Aumento de datos}
\label{sec:aumento}
Para disponer de más de 35\,000 imágenes con la misma cantidad por clase en cada partición, a partir de 5\,123 ROI originales, se recurrió a aumento de datos. El aumento se hizo siempre \emph{después} de dividir y \emph{dentro} de cada partición, para que ninguna copia cruce de partición: remuestrear antes de dividir sesga al alza las métricas \cite{demircioglu2024}. La cantidad de copias por ROI se calcula para que cada clase alcance la misma meta por partición: 8\,167 imágenes por clase en entrenamiento y 1\,751 en validación y en prueba, es decir 35\,007 en total (Tabla~\ref{tab:35k} y Fig.~\ref{fig:balance}).

\begin{table*}[t]
\centering\footnotesize
\caption{Imágenes finales por partición y clase (mismo número por clase en cada partición), ROI originales e imágenes por original}
\label{tab:35k}
\begin{tabular}{llrrrr}
\toprule
Partición & Clase & Originales & Aumentadas & Total & Imágenes por original \\
\midrule
Entrenamiento & Normal       & 771    & 7\,396 & 8\,167 & 10,6 \\
              & Degeneración & 194    & 7\,973 & 8\,167 & 42,1 \\
              & Protrusión   & 2\,612 & 5\,555 & 8\,167 & 3,1  \\
Validación    & Normal       & 170    & 1\,581 & 1\,751 & 10,3 \\
              & Degeneración & 44     & 1\,707 & 1\,751 & 39,8 \\
              & Protrusión   & 564    & 1\,187 & 1\,751 & 3,1  \\
Prueba        & Normal       & 164    & 1\,587 & 1\,751 & 10,7 \\
              & Degeneración & 42     & 1\,709 & 1\,751 & 41,7 \\
              & Protrusión   & 562    & 1\,189 & 1\,751 & 3,1  \\
\midrule
Total         &              & 5\,123 & 29\,884 & 35\,007 & \\
\bottomrule
\end{tabular}
\end{table*}

\subsubsection{Transformaciones}
Se usaron transformaciones geométricas y fotométricas leves, habituales en imagen médica \cite{islam2024}: rotación de $\pm 8^{\circ}$, desplazamiento de hasta $\pm 8\,\%$ del tamaño, zoom entre 0,90 y 1,10, y brillo y contraste de $\pm 10\,\%$, con reflejo en los bordes. No se usan reflexiones (el eje anterior--posterior tiene significado anatómico), deformación elástica ni realce. Cada copia se genera con una semilla derivada de un hash (SHA-256) del nombre del ROI, la semilla global y el índice de la copia, por lo que el proceso es determinista y reproducible.

\begin{figure}[!htb]
\centering
\begin{tikzpicture}
\begin{axis}[
  ybar, bar width=10pt, width=\columnwidth, height=4.6cm,
  symbolic x coords={Normal,Degeneración,Protrusión}, xtick=data,
  ymin=0, ymax=9800, ylabel={Imágenes (entrenamiento)},
  enlarge x limits=0.25, legend style={at={(0.5,1.03)},anchor=south,legend columns=2,font=\scriptsize},
  nodes near coords, every node near coord/.append style={font=\tiny, /pgf/number format/.cd,1000 sep={\,}},
  tick label style={font=\scriptsize}, ylabel style={font=\scriptsize},
  yticklabel style={/pgf/number format/.cd,1000 sep={\,}}, scaled y ticks=false]
\addplot coordinates {(Normal,771) (Degeneración,194) (Protrusión,2612)};
\addplot coordinates {(Normal,8167) (Degeneración,8167) (Protrusión,8167)};
\legend{Originales, Tras el aumento}
\end{axis}
\end{tikzpicture}
\caption{Balance por conteo en entrenamiento: ROI originales y total tras el aumento.}
\label{fig:balance}
\end{figure}

\subsubsection{Trazabilidad}
Cada archivo se llama \arch{T2_TSE_SAG__0188_008_L4-L5__orig.png} (original) o \arch{..._aug012.png} (copia), con paciente, corte y nivel. El manifiesto \arch{manifest_35k.csv} registra, para cada una de las 35\,007 imágenes, su partición, clase, paciente, nivel, ROI original, revisor y, en las copias, los parámetros exactos (semilla, ángulo, zoom, desplazamientos, contraste y brillo); \arch{split_pacientes.csv} guarda la partición de cada paciente. Las imágenes no se suben al repositorio: se distribuyen como un archivo comprimido de 405~MB, y el repositorio guarda los manifiestos y los scripts que las reproducen a partir de los datos públicos.

\subsubsection{Verificaciones}
El script comprueba que (i) los conjuntos de pacientes de las tres particiones son disjuntos, (ii) toda copia está en la misma partición que su original, (iii) los nombres de archivo son únicos, (iv) todo PNG está íntegro y (v) el total supera los 35\,000. Las copias respetan los rangos declarados (ángulo entre $-8^{\circ}$ y $8^{\circ}$, zoom y contraste entre 0,90 y 1,10).

% ---------------------------------------------------------------------
\subsection{Balance de clases y sesgo por nivel}
\label{sec:balance}
El balance de 1/3 por clase es \emph{por conteo}, no por originales. Entre los 5\,123 ROI originales, la Protrusión es el 73,0\,\% (3\,738), Normal el 21,6\,\% (1\,105) y Degeneración el 5,5\,\% (280); un clasificador que siempre dijera Protrusión acertaría cerca del 73\,\%. Las copias igualan la contribución de cada clase al entrenamiento, pero no añaden anatomía nueva: el tamaño efectivo del dataset es el número de pacientes (399), y la Degeneración, repetida unas 42 veces, solo tiene 42 originales en la partición de prueba, lo que deja un margen de error amplio en su recall. Para que las métricas no se inflen con copias casi idénticas, los resultados principales se calculan sobre los \textbf{originales} de la partición de prueba, y el conjunto completo se reporta aparte. Además, un modelo entrenado con prior uniforme tiende a sobrepredecir las clases raras en la población real.

La Tabla~\ref{tab:nivel} muestra un sesgo importante: casi toda la patología se concentra en L4-L5 y L5-S1. Un clasificador que solo adivinara la clase más frecuente de cada nivel acertaría 3\,842 de 5\,123 ROI (75,0\,\%), de modo que el nivel podría actuar como atajo; el aprendizaje por atajos hace que el desempeño medido sobreestime el real \cite{onglt2024}. Por ello los resultados se reportan también por nivel. Esto relativiza cualquier umbral de exactitud: un 95\,\% medido sobre estas etiquetas mide acuerdo con etiquetas de estudiantes y puede alcanzarse en parte por el atajo, y cifras internas altas suelen bajar al evaluar con datos externos \cite{yi2023}. Por eso se acompaña la exactitud con F1 macro y recall por clase: con clases desbalanceadas, las métricas habituales quedan dominadas por la clase mayoritaria \cite{mosquera2024}.

\begin{table}[!htb]
\centering\footnotesize
\setlength{\tabcolsep}{4pt}
\caption{ROI originales elegibles por nivel y clase}
\label{tab:nivel}
\begin{tabular}{lrrrrr}
\toprule
Nivel & Normal & Degen. & Protr. & Total & \% Protr. \\
\midrule
L1-L2 & 505 & 121 & 401    & 1\,027 & 39,0 \\
L2-L3 & 383 & 94  & 550    & 1\,027 & 53,6 \\
L3-L4 & 172 & 44  & 811    & 1\,027 & 79,0 \\
L4-L5 & 15  & 5   & 1\,005 & 1\,025 & 98,0 \\
L5-S1 & 30  & 16  & 971    & 1\,017 & 95,5 \\
\bottomrule
\end{tabular}
\end{table}

% III. Metodología (continuación): formulación, modelos y protocolo de evaluación
% Subsecciones H a J. Basado en la sección de selección de modelos del Hito 2
% y en docs/resultados/ (metadata.json de cada modelo).

% ---------------------------------------------------------------------
\subsection{Formulación del problema y métricas}
\label{sec:formulacion}
La tarea es de clasificación multiclase supervisada: dado un ROI $x$ de un disco, el modelo predice una clase $\hat{y}\in\{\text{Degeneración},\ \text{Normal},\ \text{Protrusión}\}$. Las clases son nominales: la prioridad usada al etiquetar (Sección~\ref{sec:anotacion}) no es un orden de gravedad, por lo que no se planteó como regresión ni como clasificación ordinal.

Las métricas son la exactitud, la precisión, el \emph{recall} y el F1 por clase y su promedio macro, la exactitud balanceada (igual al \emph{recall} macro), el área bajo la curva ROC macro (uno contra el resto) y la pérdida logarítmica. La exactitud sola es engañosa con este desbalance: un clasificador que siempre dijera Protrusión acertaría el 73\,\% de los originales de prueba. Por eso el criterio principal de comparación es el F1 macro, acompañado del \emph{recall} por clase \cite{mosquera2024}. Para responder a la pregunta P3 se reporta además la tasa de falsos negativos clínicos: la fracción de discos patológicos (Degeneración o Protrusión) que el modelo clasifica como Normal.

% ---------------------------------------------------------------------
\subsection{Modelos y ajuste}
\label{sec:modelos}
Se eligieron cuatro modelos que cumplen cuatro criterios: (i) sirven para clasificar un ROI en tres clases; (ii) tienen antecedentes en RM lumbar o en tareas médicas comparables; (iii) representan familias distintas, para que la comparación diga algo; y (iv) pueden entrenarse con la misma división por paciente y evaluarse con las mismas métricas. Todos parten de pesos preentrenados en ImageNet. Los ROI de 128$\times$128~px en escala de grises se convierten a tres canales, se redimensionan a 224$\times$224~px y se escalan a $[0,1]$; en la ResNet-50 se normalizan además con la media y la desviación estándar de ImageNet. La Tabla~\ref{tab:hiper} resume la configuración de cada modelo.

\begin{table*}[t]
\centering\footnotesize
\caption{Configuración de los cuatro modelos. Los valores elegidos entre varios se fijaron con el F1 macro (o la exactitud, en YOLO) de la validación}
\label{tab:hiper}
\begin{tabular}{>{\raggedright\arraybackslash}p{0.15\textwidth}>{\raggedright\arraybackslash}p{0.2\textwidth}>{\raggedright\arraybackslash}p{0.19\textwidth}>{\raggedright\arraybackslash}p{0.36\textwidth}}
\toprule
Modelo & Entrada & Optimización y pérdida & Hiperparámetros \\
\midrule
1. YOLOv8s-cls & Imagen 224$\times$224 & AdamW; entropía cruzada & Tasa inicial $5\times10^{-4}$, final $\times 0{,}1$, coseno con 1 época de calentamiento; decaimiento de pesos 0,01; 4 capas congeladas; \emph{dropout} 0,3; lote 64; hasta 15 épocas, paciencia 4 \\
2. ResNet-50 & Imagen 224$\times$224, normalización ImageNet & AdamW; entropía cruzada con suavizado 0,1 & Tasa $10^{-4}$ con decaimiento coseno; decaimiento de pesos 0,01; \emph{dropout} 0,3; lote 64; hasta 10 épocas, paciencia 4; mejor época: 2 \\
3. SVM & 2\,048 características de la ResNet-50 ajustada $\to$ PCA 256 & Pérdida bisagra (SVC) & Núcleo RBF, $\gamma$ = \texttt{scale}, uno contra el resto; $C\in\{0{,}1;\ 0{,}5;\ 2;\ 8\}$, elegido $C=0{,}1$ \\
4. Random Forest & Las mismas que el SVM & Índice de Gini & 300 árboles; profundidad máxima $\in\{10;\ 20;\ \text{sin límite}\}$, elegida 20 \\
\bottomrule
\end{tabular}
\end{table*}

\subsubsection{YOLOv8s-cls}
Es la variante de clasificación de la familia YOLO, la más usada en la literatura revisada \cite{liawrungrueang2023,prisilla2023,wang2025}, y es coherente con el resto del flujo, que ya usa YOLO para detectar los discos. Se ajustó con AdamW, entropía cruzada, las cuatro primeras capas congeladas y \emph{dropout} de 0,3. Como el aumento principal ya está en los archivos, el aumento en línea se limitó a variar la proporción del recorte (de 3:4 a 4:3) y el brillo ($\pm 5\,\%$), sin reflexiones.

\subsubsection{ResNet-50}
Es una red convolucional de otra familia (residual), con respaldo en la clasificación de discos lumbares \cite{tan2024,alkubaisi2022}. Se congelaron la primera convolución y el primer bloque residual, y la capa final se sustituyó por \emph{dropout} de 0,3 y una capa lineal de tres salidas. Se entrenó con AdamW, tasa de aprendizaje con decaimiento coseno, entropía cruzada con suavizado de etiquetas de 0,1 y un aumento en línea suave (rotación de $\pm 6^{\circ}$, desplazamiento del 4\,\%, zoom de 0,95 a 1,05, y brillo y contraste de $\pm 10\,\%$). La mejor época se eligió por F1 macro en validación, con parada temprana.

\subsubsection{SVM y Random Forest sobre características profundas}
Un clasificador clásico sobre los píxeles crudos trata cada valor por separado y no reconoce la forma. Por eso se usó la ResNet-50 como extractor: sin su capa final, entrega 2\,048 valores por imagen, que se estandarizan, se reducen a 256 componentes con análisis de componentes principales (PCA; 80,0\,\% de la varianza) y se clasifican con un SVM de núcleo RBF o con un Random Forest de 300 árboles. La técnica tiene antecedentes sobre este mismo corpus \cite{sari2025} y en otras tareas médicas \cite{khoiwal2024}; el Random Forest solo cuenta con antecedentes en datos tabulares \cite{zhang2023} y se presenta como una extensión propia. El escalador, el PCA y los clasificadores se ajustan solo con el entrenamiento.

El diseño inicial usaba la ResNet-50 con los pesos de ImageNet congelados, para que el SVM y el Random Forest no dependieran de la red ajustada. En validación rindieron poco (F1 macro de 0,64 y 0,60), así que se adoptó como extractor la ResNet-50 ya ajustada (F1 de 0,80 y 0,79). El costo es que los modelos 3 y 4 dependen del modelo 2. La versión con ImageNet se conserva como ablación (Sección~\ref{sec:res_p2}). Como líneas base del curso se entrenaron también una regresión logística y un KNN sobre las mismas características, que no cuentan entre los cuatro modelos.

% ---------------------------------------------------------------------
\subsection{Protocolo de evaluación}
\label{sec:protocolo}
Los cuatro modelos usan la misma división por paciente (Tabla~\ref{tab:split}) y la misma semilla (42). Las épocas, los hiperparámetros y las variantes se eligieron con la validación; la prueba se evaluó una sola vez, con la configuración final. Las variantes probadas en validación (otras semillas, tasas de aprendizaje, suavizado y ensambles) quedan registradas en el repositorio.

La guía del curso recomienda aumentar solo el entrenamiento. Aquí también se aumentaron la validación y la prueba, para tener la misma cantidad por clase (Sección~\ref{sec:aumento}). Las copias de validación y prueba nunca entran al entrenamiento, por lo que no hay fuga; aun así, para que las copias casi idénticas no inflen las métricas, \textbf{los resultados principales se calculan sobre los 768 ROI originales de la prueba}, con la proporción natural de clases (42 de Degeneración, 164 Normales y 562 de Protrusión). El conjunto de prueba completo, balanceado y con copias (5\,253 imágenes), se reporta aparte. Con 768 imágenes, el intervalo de confianza del 95\,\% de la exactitud es de unos $\pm 2{,}2$ puntos (intervalo de Wilson).


\section{Resultados y discusión}
\label{sec:resultados}
% IV. Resultados y discusión
% Cifras: docs/resultados/comparacion/*.csv y matrices confusion_test_orig.csv
Salvo que se indique lo contrario, las cifras corresponden a los 768 ROI originales de la partición de prueba (60 pacientes), evaluada una sola vez.

% ---------------------------------------------------------------------
\subsection{Distribución de clases y diferencias visuales (P1)}
\label{sec:res_p1}
Tras la revisión humana, la Protrusión es el 73,0\,\% de los 5\,123 ROI originales, Normal el 21,6\,\% y Degeneración el 5,5\,\% (Sección~\ref{sec:balance}). La distribución depende mucho del nivel: la Protrusión pasa del 39,0\,\% en L1-L2 al 98,0\,\% en L4-L5 (Tabla~\ref{tab:nivel}). La Degeneración, en cambio, aparece sobre todo en L1-L2 y L2-L3, en parte porque en los niveles inferiores un disco degenerado que además protruye recibe la etiqueta Protrusión.

Las diferencias visuales que guiaron el etiquetado son las de la escala de Pfirrmann y la nomenclatura de consenso \cite{pfirrmann2001,fardon2014}. Un disco Normal tiene el núcleo brillante en T2, bien distinguible del anillo, y altura conservada. En la Degeneración el disco se oscurece por pérdida de hidratación y pierde altura, sin que el material discal sobrepase el contorno vertebral. En la Protrusión el material discal se extiende más allá de ese contorno, casi siempre hacia el canal (Fig.~\ref{fig:rois}). El brillo distingue, por tanto, Normal de Degeneración, y la forma del contorno posterior distingue la Protrusión. La separación entre abombamiento y hernia depende de la extensión sobre la circunferencia del disco, que no se ve en un corte sagital, y por eso esas dos clases se fusionaron. Las dos patologías, además, coexisten con frecuencia: un disco degenerado puede protruir, y en ese caso recibe la etiqueta Protrusión por la regla de prioridad, lo que anticipa la confusión entre ambas clases.

% ---------------------------------------------------------------------
\subsection{Comparación de los cuatro modelos}
\label{sec:res_comparacion}
La Tabla~\ref{tab:resultados} resume las métricas. La ResNet-50 obtuvo la mayor exactitud (88,9\,\%, 683 de 768; IC 95\,\%: 86,5--91,0) y el mayor F1 macro (0,769), seguida por el SVM y el Random Forest (88,4\,\%) y el YOLOv8s-cls (87,8\,\%). Los intervalos de confianza de los cuatro se solapan: las diferencias de exactitud (un punto como máximo) no son concluyentes. El Random Forest obtuvo el mayor ROC-AUC (0,942) y la menor pérdida logarítmica (0,327), es decir, las probabilidades mejor ordenadas y calibradas, aunque su F1 macro es el más bajo de los cuatro. El ensamble de los cuatro modelos no mejora a la ResNet-50 en exactitud ni en F1.

\begin{table*}[t]
\centering\footnotesize
\setlength{\tabcolsep}{4.5pt}
\caption{Resultados en la partición de prueba. Las columnas de originales usan los 768 ROI originales; las métricas son promedios macro. ``Entren.'' es la exactitud en los originales de entrenamiento y ``Prueba completa'' el conjunto balanceado con copias (5\,253 imágenes). En negrita, el mejor de los cuatro modelos}
\label{tab:resultados}
\begin{tabular}{lrrrrrrrrr}
\toprule
 & \multicolumn{6}{c}{Prueba, originales (768)} & & \multicolumn{2}{c}{Prueba completa} \\
\cmidrule(lr){2-7}\cmidrule(lr){9-10}
Modelo & Exactitud & Precisión & \emph{Recall} & F1 & ROC-AUC & Pérd. log. & Entren. & Exactitud & F1 \\
\midrule
1. YOLOv8s-cls        & 0,878 & 0,771 & 0,732 & 0,750 & 0,931 & 0,488 & 0,984 & 0,759 & 0,752 \\
2. ResNet-50          & \textbf{0,889} & 0,795 & \textbf{0,751} & \textbf{0,769} & 0,919 & 0,329 & 0,971 & \textbf{0,775} & \textbf{0,768} \\
3. SVM (RBF)          & 0,884 & 0,797 & 0,718 & 0,748 & 0,899 & -- & 0,979 & 0,754 & 0,746 \\
4. Random Forest      & 0,884 & \textbf{0,810} & 0,710 & 0,746 & \textbf{0,942} & \textbf{0,327} & 0,998 & 0,743 & 0,731 \\
\midrule
\multicolumn{10}{l}{\emph{Comparaciones adicionales}} \\
Ensamble de los 4 (promedio)            & 0,889 & 0,782 & 0,721 & 0,745 & 0,947 & 0,321 & 0,989 & 0,762 & 0,752 \\
Regresión logística (línea base)        & 0,878 & 0,819 & 0,696 & 0,740 & 0,896 & 0,732 & 0,988 & 0,730 & 0,717 \\
KNN (línea base)                        & 0,888 & 0,794 & 0,722 & 0,751 & 0,858 & 1,000 & 0,980 & 0,758 & 0,749 \\
SVM, características de ImageNet        & 0,776 & 0,563 & 0,600 & 0,576 & 0,829 & -- & 0,841 & 0,633 & 0,619 \\
Random Forest, características de ImageNet & 0,753 & 0,518 & 0,546 & 0,524 & 0,831 & 0,726 & 0,999 & 0,598 & 0,577 \\
\bottomrule
\multicolumn{10}{l}{\scriptsize El SVM no entrega probabilidades, por lo que no tiene pérdida logarítmica.}
\end{tabular}
\end{table*}

En el conjunto de prueba completo, balanceado, la exactitud baja a 74--78\,\%. Ese conjunto pesa igual las tres clases, de modo que la Degeneración, la clase más difícil, cuenta tanto como la Protrusión. La diferencia entre las dos cifras muestra cuánto de la exactitud sobre originales proviene de la clase mayoritaria.

% ---------------------------------------------------------------------
\subsection{Resultados por clase y falsos negativos (P3)}
\label{sec:res_p3}
La Tabla~\ref{tab:porclase} y la Fig.~\ref{fig:confusion} muestran el comportamiento por clase. Los cuatro modelos reconocen la Protrusión con un \emph{recall} de 0,94 a 0,95 y la clase Normal con 0,76 a 0,81. La Degeneración es la clase débil: la ResNet-50 y el YOLOv8s-cls detectan la mitad de los 42 casos (IC 95\,\%: 0,36--0,64), y el SVM y el Random Forest, entre el 40 y el 43\,\%. Sus errores van sobre todo hacia Protrusión (de 12 a 17 casos según el modelo), lo que concuerda con la coexistencia de ambas patologías descrita en la Sección~\ref{sec:res_p1}. La segunda confusión frecuente es Normal predicho como Protrusión (de 26 a 36 casos).

\begin{figure*}[t]
\centering
\includegraphics[width=0.95\textwidth]{matrices_confusion_es.png}
\caption{Matrices de confusión de los cuatro modelos en los 768 ROI originales de prueba. Cada celda muestra el número de casos y el porcentaje de su fila (clase real).}
\label{fig:confusion}
\end{figure*}

Si se agrupan Degeneración y Protrusión como ``patológico'', los cuatro modelos clasifican como Normal entre el 5,0 y el 5,5\,\% de los 604 discos patológicos; es decir, su sensibilidad para detectar patología es del 94,5 al 95,0\,\%. En falsos negativos clínicos, el SVM y el Random Forest son ligeramente mejores (30 casos frente a 32 y 33), pero a costa de un menor \emph{recall} en Degeneración. Ningún modelo domina en todas las métricas.

\begin{table}[!htb]
\centering\footnotesize
\setlength{\tabcolsep}{3pt}
\caption{Recall y F1 por clase en los originales de prueba, y fracción de discos patológicos (604) clasificados como Normal (FN)}
\label{tab:porclase}
\begin{tabular}{lrrrrrrr}
\toprule
 & \multicolumn{3}{c}{\emph{Recall}} & \multicolumn{3}{c}{F1} & \\
\cmidrule(lr){2-4}\cmidrule(lr){5-7}
Modelo & Degen. & Normal & Protr. & Degen. & Normal & Protr. & FN \\
\midrule
YOLOv8s-cls   & 0,500 & 0,756 & 0,941 & 0,545 & 0,775 & 0,929 & 5,3\,\% \\
ResNet-50     & 0,500 & 0,811 & 0,941 & 0,568 & 0,806 & 0,935 & 5,5\,\% \\
SVM           & 0,429 & 0,774 & 0,950 & 0,522 & 0,791 & 0,932 & 5,0\,\% \\
Random Forest & 0,405 & 0,774 & 0,952 & 0,515 & 0,791 & 0,931 & 5,0\,\% \\
\bottomrule
\end{tabular}
\end{table}

Respecto a la pregunta P3, la ResNet-50 es la arquitectura con mejor exactitud y F1 macro, pero ninguna alcanza la meta del 90\,\%: la ResNet-50 queda a 1,1 puntos, dentro de su intervalo de confianza, pero por debajo. En validación, todas las variantes probadas (dos arquitecturas, tres semillas, regularización y ensambles de hasta ocho redes) se estancaron entre el 86 y el 87,5\,\% sobre los originales. Que métodos tan distintos choquen con el mismo techo indica que el límite está en los datos más que en el modelo.

% ---------------------------------------------------------------------
\subsection{Clasificadores clásicos frente a redes profundas (P2)}
\label{sec:res_p2}
Sobre características de la ResNet-50 ajustada, el SVM y el Random Forest igualan a las redes (88,4\,\% de exactitud y F1 macro de 0,748 y 0,746), y lo mismo ocurre con la regresión logística y el KNN (Tabla~\ref{tab:resultados}). Con características de ImageNet sin ajustar, en cambio, caen al 77,6 y 75,3\,\% de exactitud y a un F1 macro de 0,576 y 0,524, y la Degeneración casi desaparece (\emph{recall} de 0,21 y 0,10). La respuesta a P2 es, por tanto, matizada: un clasificador clásico puede ser tan confiable como una red profunda, pero solo si trabaja sobre características que ya aprendió una red ajustada a la tarea. En ese caso, lo que distingue las clases está en la representación, no en el clasificador final. No se probaron características diseñadas a mano (intensidad media del disco, altura o contorno), que quedan como trabajo futuro.

% ---------------------------------------------------------------------
\subsection{Sobreajuste}
\label{sec:res_sobreajuste}
Los cuatro modelos aciertan entre el 97,1 y el 99,8\,\% de los originales de entrenamiento, frente al 87,8--88,9\,\% de la prueba (Tabla~\ref{tab:resultados}). Las redes alcanzan su mejor época en validación muy pronto (la ResNet-50 en la época 2), y a partir de ahí el error de entrenamiento sigue bajando sin que mejore la validación. El Random Forest es el caso extremo (99,8\,\% en entrenamiento). Las medidas de regularización (capas congeladas, \emph{dropout}, decaimiento de pesos, suavizado de etiquetas, aumento en línea y parada temprana) contuvieron el sobreajuste sin eliminarlo. Es lo esperable con 279 pacientes de entrenamiento y 194 originales de Degeneración repetidos unas 42 veces cada uno.

% ---------------------------------------------------------------------
\subsection{Sesgo por nivel vertebral}
\label{sec:res_nivel}
La Tabla~\ref{tab:nivelres} compara la exactitud de la ResNet-50 en cada nivel con la de un clasificador trivial que siempre respondiera la clase más frecuente de ese nivel en la prueba. El trivial acierta el 74,7\,\% del total, muy cerca del 73\,\% de responder siempre Protrusión. En L4-L5 y L5-S1, donde casi todo es Protrusión, la ResNet-50 apenas supera al trivial (97,4 frente a 94,8\,\% y 94,2 frente a 94,2\,\%). Su ventaja real está en L1-L2 y L2-L3, donde las clases están más repartidas (de 30 puntos en L1-L2 y 29 en L2-L3). Así, una parte de la exactitud global puede venir del nivel y no de la imagen, como ocurre con el aprendizaje por atajos \cite{onglt2024}; los niveles inferiores, donde el atajo es más fácil, pesan lo mismo en la exactitud global que los superiores.

\begin{table}[!htb]
\centering\footnotesize
\caption{Exactitud por nivel de la ResNet-50 frente a un clasificador que responde la clase más frecuente de cada nivel (originales de prueba)}
\label{tab:nivelres}
\begin{tabular}{lrrr}
\toprule
Nivel & ROI & Clase más frecuente & ResNet-50 \\
\midrule
L1-L2 & 153 & 48,4\,\% & 78,4\,\% \\
L2-L3 & 153 & 56,2\,\% & 85,6\,\% \\
L3-L4 & 154 & 79,9\,\% & 89,0\,\% \\
L4-L5 & 154 & 94,8\,\% & 97,4\,\% \\
L5-S1 & 154 & 94,2\,\% & 94,2\,\% \\
\midrule
Total & 768 & 74,7\,\% & 88,9\,\% \\
\bottomrule
\end{tabular}
\end{table}

% ---------------------------------------------------------------------
\subsection{Comparación con el estado del arte}
\label{sec:res_sota}
Las cifras de este trabajo son inferiores a las de la literatura (97--99\,\%) \cite{sari2025,liawrungrueang2023}, pero no son comparables de forma directa. Sarı y Atasoy usaron el mismo corpus, pero en una tarea binaria, con cortes T1 y validación cruzada sin división declarada por paciente \cite{sari2025}; aquí la tarea tiene tres clases muy desbalanceadas, la división es por paciente y la evaluación se hace sobre originales que el modelo nunca vio. Natalia \emph{et al.}, sobre el mismo corpus y prediciendo el grado de Pfirrmann, reportaron una exactitud media del 88,1\,\% \cite{natalia2024}, del mismo orden que la obtenida aquí. Además, las cifras internas altas suelen bajar en validación externa \cite{yi2023,nigru2024}, que este trabajo no pudo hacer. Por último, las etiquetas son de estudiantes, sin radiólogo, de modo que la exactitud mide el acuerdo con esas etiquetas y no con un diagnóstico clínico \cite{shi2024}.


\section{Conclusiones}
\label{sec:conclusiones}
% V. Conclusiones
Este trabajo construyó, a partir de un corpus público de 515 pacientes, un dataset trazable de 5\,123 ROI de discos lumbares de 399 pacientes en tres clases, y comparó sobre él cuatro modelos con la misma división por paciente y las mismas métricas. Las respuestas a las preguntas de investigación son las siguientes.

\textbf{P1.} El dataset está muy desbalanceado: Protrusión es el 73,0\,\% de los ROI, Normal el 21,6\,\% y Degeneración el 5,5\,\%, y la patología se concentra en L4-L5 y L5-S1. La señal T2 del disco separa Normal de Degeneración y la forma del contorno posterior identifica la Protrusión; Hernia y Abombamiento no pueden separarse en el plano sagital, lo que obligó a fusionarlos.

\textbf{P2.} El SVM y el Random Forest alcanzan a las redes (88,4\,\% de exactitud) solo cuando usan características de una ResNet-50 ajustada a la tarea; con características genéricas de ImageNet caen al 75--78\,\% y casi no detectan la Degeneración. Son más rápidos de ajustar, pero no sustituyen a la red: dependen de ella.

\textbf{P3.} La ResNet-50 obtuvo la mayor exactitud (88,9\,\%) y el mayor F1 macro (0,769), sin diferencias concluyentes con los otros tres modelos y sin alcanzar la meta del 90\,\%. Los cuatro clasifican como Normal solo entre el 5,0 y el 5,5\,\% de los discos patológicos, pero detectan apenas entre el 40 y el 50\,\% de los casos de Degeneración. Una parte de la exactitud se explica por el nivel vertebral.

En conjunto, el factor limitante no fue la arquitectura sino los datos: pocas imágenes originales de Degeneración, etiquetas de estudiantes sin radiólogo y un sesgo fuerte por nivel. Por eso las métricas por clase y por nivel son tan importantes como la exactitud global.


\subsection{Limitaciones y trabajo futuro}
% Limitaciones y trabajo futuro. Va dentro de "V. Conclusiones".
Las principales limitaciones son las siguientes:
\begin{itemize}
  \item Las etiquetas son de estudiantes, sin radiólogo, revisadas viendo la sugerencia de un modelo, y no se midió el acuerdo entre anotadores (Sección~\ref{sec:anotacion}).
  \item Solo hay plano sagital; Hernia y Abombamiento no son separables y se fusionaron después de ver los datos.
  \item Un único equipo de resonancia de 1,5~T, sin diversidad de escáner ni validación externa, que la literatura recomienda para estos modelos \cite{nigru2024,wang2025}.
  \item El tamaño efectivo son 399 pacientes; las copias no aportan anatomía nueva. Degeneración cuenta con pocos pacientes (81, 17 y 22 en las tres particiones) y solo 42 originales en la prueba, lo que deja un margen de error amplio en su \emph{recall}.
  \item Sesgo por nivel y prior uniforme (Secciones~\ref{sec:balance} y \ref{sec:res_nivel}).
  \item El conjunto público no incluye un identificador de paciente: el campo \texttt{PatientID} está vacío en todos los archivos \texttt{.IMA} originales, algo que depende por completo de la fuente y no de este trabajo. Por eso se usó el número de carpeta como identificador de paciente, siguiendo la organización de la fuente (una carpeta por paciente, 515 en total), y la división por paciente se hizo por carpeta. Entre las 414 carpetas usadas no se repite ningún \texttt{StudyInstanceUID} ni ningún \texttt{StudyTime}, así que no hay indicios de duplicados; solo si la propia fuente hubiera repartido a una misma persona en dos carpetas podría existir un solapamiento entre particiones, y eso no puede verificarse sin el identificador original.
  \item Un error en la selección de cortes hizo perder 36 pacientes y dejó a otros 8 con un corte parasagital (Sección~\ref{sec:depuracion}).
  \item Las métricas del detector provienen de una repetición del entrenamiento y no del modelo que generó los ROI.
  \item El SVM y el Random Forest dependen de la ResNet-50 ajustada, por lo que no son modelos independientes, y cada modelo oficial corresponde a una sola semilla.
  \item El sistema es una herramienta académica, no diagnóstica.
\end{itemize}

Como trabajo futuro, las etiquetas de Degeneración podrían contrastarse con los grados de Pfirrmann validados por radiólogo que existen para este mismo corpus en L3-L4, L4-L5 y L5-S1 \cite{masood2026}, y los modelos podrían evaluarse en conjuntos multicéntricos anotados por expertos \cite{richards2026}. Corregir el filtrado de cortes, seleccionando solo la serie T2 TSE, recuperaría además los 36 pacientes perdidos. Como Degeneración y Protrusión coexisten en un mismo disco, una formulación multietiqueta, en lugar de la regla de prioridad, podría reducir la confusión entre ambas. Por último, incorporar los cortes axiales permitiría separar de nuevo Hernia y Abombamiento, y unas características diseñadas a mano (intensidad y altura del disco) completarían la respuesta a la pregunta P2.


\section*{Reconocimientos}
Los autores agradecen a S. Sudirman y a sus coautores por publicar en acceso abierto el conjunto de datos que hizo posible este trabajo \cite{sudirman2019}.

\bibliographystyle{IEEEtran}
\bibliography{referencias}

\end{document}
